% !TeX program = lualatex
% !TeX encoding = utf8
% !TeX root = ../Thermod.tex

\setcounter{chapter}{12}

%=========================================================
\Opensolutionfile{answer}[\currfilebase/\currfilebase-Answers]
\Writetofile{answer}{\protect\section*{\nameref*{\currfilebase}}}
\chapter{Статистичні розподіли}\label{\currfilebase}
\graphicspath{{\currfilebase/pictures}}
\makeatletter
\def\input@path{{\currfilebase/}}
\makeatother

%% --------------------------------------------------------
\section{128}
%% --------------------------------------------------------

Визначаючи тиск газу, який діє на стінку судини, ми задовольнялися інтуїтивними уявленнями про рух молекул газу, що має хаотичний характер, анітрохи не замислюючись про спроби його опису. Досить було зробити розумне припущення щодо ізотропності цього руху і, як наслідок, незалежності компонент швидкості. Саме тому можемо стверджувати, що $\langle v_x v_y \rangle = 0$, $\langle v_x^2 \rangle = \langle v_y^2 \rangle = \langle v_z^2 \rangle = \frac{1}{3} \langle v^2 \rangle$ і т.д.

Припущення це видається прийнятним, оскільки в протилежному випадку рух молекул газу носив би спрямований характер, що важко очікувати у відсутності зовнішніх полів та інших сил. Надалі ми будемо дотримуватися цього припущення, проте спробуємо знайти взаємозв’язок між характеристиками руху і якими-небудь параметрами, котрі характеризують термодинамічну систему. В якійсь мірі це завдання нами було вирішене, оскільки знайдено, що температура, а також $pV$ пропорційні середньому квадрату швидкості руху молекул.

Для того щоб з'ясувати, яким же чином можна охарактеризувати рух молекул, розберемося детально, що являє собою процедура усереднення. За визначенням середнього, якщо будь-яка величина, нехай це буде швидкість $v$, набуває в серії експериментів $N$ різних значень, то її середнє значення може бути знайдено як:
%
\begin{equation}\label{eq:128_1}
    \langle \vect{v} \rangle = \frac{1}{N} \sum_{i=1}^{N} \vect{v}_i.
    \tag{1}
\end{equation}
%
Отже, середню швидкість молекул газу, в принципі, можна визначити як середнє арифметичне значення швидкостей усіх молекул системи. Для цього необхідно було б експериментально виміряти швидкість кожної молекули, після чого обчислити суму отриманих значень і поділити на їхню кількість. Проте реалізація такого підходу на практиці є неможливою через колосальну кількість частинок у макроскопічному об’ємі газу.

Виділимо, однак, у всьому масиві молекул ті з них, швидкості яких близькі одна до одної, скажімо, лежать в діапазоні $1$ см/сек. Тоді, з високою точністю \eqref{eq:128_1} можна переписати як:
%
\begin{equation}\label{eq:128_2}
    \langle \vect{v} \rangle = \sum_{k=1}^{M} \frac{N_k}N\vect{v}_k = \sum_{k=1}^{M} \vect{v}_k \cdot \Delta \mathcal{P}(k),
    \tag{2}
\end{equation}
%
де $M$ --- кількість інтервалів, а $N_k$ --- кількість молекул, швидкості яких, з точністю $\pm 0.5$ см/сек рівні $\vec{v_k}$. Величина $\Delta P(k) = N_k / N$ має сенс частоти, з якою довільно обрана молекула буде знайдена зі швидкістю $v_k$. При $N \to \infty$ ці величини набувають властивість ймовірності знаходження в статистичному ансамблі швидкості $v_k$.

Оскільки швидкість --- векторна величина, то зручно ввести простір швидкостей $v_x, v_y, v_z$. Кінці векторів, відкладених з початку координат, будуть характеризувати швидкості молекул. Елементарний об’єм в цьому просторі буде $dv_x dv_y dv_z$. Перетворимо \eqref{eq:128_2} до вигляду:
%
\begin{equation}\label{eq:128_3}
    \langle v \rangle \approx \sum_{k=1}^{M} v_k f(v_k) \Delta v,
\end{equation}
%
де припускаємо, що існує границя $f(v) = \lim_{\Delta v \to 0} \Delta P / \Delta v$, яка називається щільністю ймовірності знайти молекулу зі швидкістю $v$. Зробивши інтервал $\Delta v$ як завгодно малим і перейшовши в (\textit{3}) до інтегрування, знайдемо:
%
\begin{equation}
    \langle v \rangle = \int v f(v)\, dv,
    \tag{4}
\end{equation}
%
де $dv$ означає ймовірність того, що кінець вектора $v$ потрапляє в елемент об’єму $dv$, що знаходиться на відстані $v$ і в напрямку $v$ від початку координат. В силу того, що $f(v)$, як і будь-яка функція від $v$, має ту ж імовірність, можна записати:
%
\begin{equation}
    f(v_x, v_y, v_z) = f(v),
    \tag{5}
\end{equation}
%
що робить функцію $f(v)$ універсальною характеристикою мікроскопічного руху молекул.

%% --------------------------------------------------------
\section{129}
%% --------------------------------------------------------

Отже, працювати ми будемо в просторі швидкостей і шукати $f(v)$. Для початку зауважимо, що в силу ізотропності, щільність розподілу залежить тільки від значення швидкості, а не від її напрямку. В силу тих же причин ймовірність $f(v_x) dv_x$ того, що компонент швидкості $v_x$ лежить в інтервалі $[v_x, v_x + dv_x]$, не залежить від ймовірності для двох інших компонент швидкості ($f(v_y) dv_y$ та $f(v_z) dv_z$) потрапити в інтервали $[v_y, v_y + dv_y], [v_z, v_z + dv_z]$ відповідно. Отже, вид функції $f$ буде однаковий для всіх компонент. Тоді:
%
\begin{equation}
    f(v_x, v_y, v_z) = f(v_x) f(v_y) f(v_z).
    \tag{1}
\end{equation}
%
Або:
%
\begin{equation}\label{eq:2}
    f(v_x, v_y, v_z) = f(v), \quad v = \sqrt{v_x^2 + v_y^2 + v_z^2}.
    \tag{2}
\end{equation}
%
Візьмемо похідну по $v_x$ в \eqref{eq:2}. Тоді неважко переконатися, що:
%
\begin{equation}
    \frac{d}{dv_x} f(v_x, v_y, v_z) = f'(v_x)f(v_y)f(v_z).
    \tag{3}
\end{equation}
%
Або:
%
\begin{equation}
    \frac{df}{dv_x} = f'(v_x) f(v_y) f(v_z).
    \tag{4}
\end{equation}
%
Аналогічним чином, можна отримати, що $\Phi(\upsilon) = \eta(\upsilon_y) = \xi(\upsilon_z)$. Тоді рівність $\varphi(\upsilon_x) = \eta(\upsilon_y) = \xi(\upsilon_z)$ має виконуватися для будь-яких значень $\upsilon_x, \upsilon_y, \upsilon_z$, що можливо лише, якщо їх значення дорівнює деякій константі $\varphi(\upsilon_x) = \eta(\upsilon_y) = \xi(\upsilon_z) = 2\gamma$. Тоді:
%
\begin{equation}
    \ln f(\upsilon_x) = -\gamma \upsilon_x^2 + \ln a.
    \tag{5}
\end{equation}
%
Або:
%
\begin{equation}
    f(\upsilon_x) = a e^{-\gamma \upsilon_x^2}.
    \tag{6}
\end{equation}
%
Отриманий вираз симетричний відносно початку координат. Це говорить про те, що компоненти швидкості $\upsilon_x$ молекули з однаковою ймовірністю приймають додатні і від’ємні значення. Те саме можна сказати і про $\upsilon_y, \upsilon_z$.

Значення $a$ можна знайти з умови того, що $\upsilon_x$ зі 100\% ймовірністю приймає хоч якесь значення між $-\infty$ і $+\infty$:
%
\begin{equation}
    \int_{-\infty}^{\infty} f(\upsilon_x)\, d\upsilon_x = 1.
    \tag{7}
\end{equation}
%
Цю умову називають умовою нормування. Звідси знайдемо:
%
\begin{equation}
    a = \sqrt{\frac{\gamma}{\pi}}.
    \tag{8}
\end{equation}
%
Для знаходження $\gamma$ скористаємося тим, що на одну ступінь свободи поступального руху молекули в середньому припадає кінетична енергія $\langle \varepsilon_x \rangle = \frac{1}{2} kT$. Тоді:
%
\begin{equation}
    \langle \upsilon_x^2 \rangle = \frac{kT}{m}.
    \tag{9}
\end{equation}
%
Кінцевий результат такий:
%
\begin{equation}
    f(\upsilon_x) = \sqrt{\frac{m}{2\pi kT}} e^{-\frac{m \upsilon_x^2}{2kT}}.
    \tag{10}
\end{equation}
%

%% --------------------------------------------------------
\section{130}
%% --------------------------------------------------------

У більшості випадків нас цікавить тільки модуль швидкості, а не її напрямок (що дійсно так в силу ізотропності простору швидкостей). Нам слід знайти ймовірність $\varphi(\upsilon)d\upsilon$, того, що вектор швидкості приймаючи довільний напрямок, буде лежати в інтервалі $[\upsilon, \upsilon + d\upsilon]$. Це означає, що кінець вектора потрапляє в шар товщиною $d\upsilon$ на сфері радіуса $\upsilon$. Тобто в елемент об’єму $dv_x dv_y dv_z = 4\pi \upsilon^2 d\upsilon$. Тоді ймовірність потрапити в цей шар:

\begin{equation}
    \varphi(\upsilon)d\upsilon = f(\upsilon_x)f(\upsilon_y)f(\upsilon_z)\cdot 4\pi \upsilon^2 d\upsilon.
    \label{eq:probability_shell}
\end{equation}

Звідси шукана функція:

\begin{equation}
    \varphi(\upsilon) = 4\pi \left( \frac{m}{2\pi kT} \right)^{3/2} \upsilon^2 e^{-\frac{m\upsilon^2}{2kT}}.
    \label{eq:maxwell_speed}
\end{equation}

Це є вичерпною характеристикою руху молекул. Зокрема, за допомогою цієї функції розподілу можна знайти:

\begin{itemize}
 \item  Середню швидкість:
  \begin{equation}
      \langle \upsilon \rangle = \int_0^\infty \upsilon \varphi(\upsilon)\, d\upsilon.
      \label{eq:avg_speed}
  \end{equation}
\item Середньоквадратичну швидкість:
  \begin{equation}
      \upsilon_{\text{rms}} = \sqrt{\langle \upsilon^2 \rangle}.
      \label{eq:rms_speed}
  \end{equation}
\item  Найбільш ймовірну швидкість:
  \begin{equation}
      \upsilon_p = \arg\max_\upsilon \varphi(\upsilon).
      \label{eq:most_probable_speed}
  \end{equation}

\end{itemize}

Кількість молекул, що мають швидкості, які належать діапазону $[\upsilon, \upsilon + d\upsilon]$, дорівнює добутку загальної кількості частинок $N$ на ймовірність знайти такі молекули:

\begin{equation}
    dN = N \cdot \varphi(\upsilon)d\upsilon.
    \label{eq:distribution_N}
\end{equation}

Аналогічний вираз має місце і для концентрації частинок.

%% --------------------------------------------------------
\section{131}
%% --------------------------------------------------------

У випадку, якщо постановка задачі вимагає врахування відмінностей напрямів руху частинок, слід використати функцію розподілу $F(\vec{\upsilon})$, яка дорівнює \eqref{eq:maxwell_component}.

Визначимо кількість частинок $Z$, які втрапляють за одиницю часу $dt$ в одиницю площі $d\sigma$. Для цього виділимо косокутний циліндр з основою $d\sigma$ і висотою, що дорівнює $\upsilon_x dt$. Кількість молекул в цьому циліндрі:

\begin{equation}
    dz = n(\upsilon_x) \upsilon_x dt dx,
    \label{eq:cylinder_density}
\end{equation}
де $n(\upsilon_x)$ --- концентрація всіх молекул з компонентами $\upsilon_x$.

Тепер, щоб знайти кількість всіх частинок, слід просумувати по всім компонентам $\upsilon_x$, які лежать в діапазоні $[0,+\infty]$:

\begin{equation}
    Z = \int_0^\infty \int_{-\infty}^{\infty} \int_{-\infty}^{\infty} \upsilon_x f(\upsilon_x)f(\upsilon_y)f(\upsilon_z)\, d\upsilon_y d\upsilon_z d\upsilon_x.
    \label{eq:total_particles}
\end{equation}

Перші два інтеграли дорівнюють 1, в силу нормування $f(\upsilon_y)$ та $f(\upsilon_z)$. Останній інтеграл дорівнює $\frac{1}{2}\langle \upsilon \rangle$. Тоді шукана кількість частинок, яка проходить в одиницю часу через одиницю площі, дорівнює:

\begin{equation}
    Z = \frac{1}{4} n \langle \upsilon \rangle.
    \label{eq:flux_formula}
\end{equation}

Зауважимо, що спрощена оцінка на зразок обчислення кількості частинок, що рухаються до стінки, використана нами при виведенні тиску ідеального газу дає:

\begin{equation}
    Z = \frac{1}{6} n \langle \upsilon \rangle.
    \label{eq:simplified_flux}
\end{equation}

Що не значним чином відрізняється від \eqref{eq:flux_formula}.

За допомогою \eqref{eq:flux_formula} можна знайти кількість молекул, що вилітають з поверхні рідини або твердого тіла в одиницю часу. При цьому неважливо, чи знаходиться речовина в рівновазі з насиченим паром, або випаровування відбувається в вакуумі, тому що кількість випарених частинок залежить від стану тіла, а не середовища, з яким воно межує. На цьому заснований метод вимірювання тиску насиченої пари тугоплавких металів, що дорівнює:

\begin{equation}
    p = m Z = \frac{1}{4} m n \langle \upsilon \rangle.
    \label{eq:vapor_pressure}
\end{equation}

якщо прийняти розподіл молекул за швидкостями за максвеллівський.

%% --------------------------------------------------------
\section{132}
%% --------------------------------------------------------

Формула (\ref{eq:flux_formula}) може бути отримана і без залучення поняття максвеллівського розподілу. Справді, при її виведенні використовувалося тільки властивість ізотропності розподілу. Це дає надію на те, що (\ref{eq:flux_formula}) може бути отримана лише з використанням властивостей ізотропності руху газу.

В силу цієї властивості концентрація молекул, напрямки швидкостей яких утворюють з нормаллю до площі $d\sigma$ кути в межах тілесного кута $[\Omega, \Omega + d\Omega]$, становить:
%
\begin{equation}
    dn = n \cdot \frac{d\Omega}{4\pi}.
    \label{eq:solid_angle_density}
\end{equation}

Число молекул, що вдаряються в одиничну площадку за одиницю часу і мають напрямок швидкостей, що потрапляють в тілесний кут $d\Omega$, становить:
%
\begin{equation}
    dZ = v \cos\theta \cdot dn.
    \label{eq:particles_impact}
\end{equation}

Тоді повне число ударів становить:
%
\begin{equation}
    Z = \int_\Omega v \cos\theta \cdot \frac{n}{4\pi} d\Omega.
    \label{eq:total_impacts}
\end{equation}

При обчисленні інтеграла враховано, що $\int \cos\theta \, d\Omega = 2\pi$. Отже:
%
\begin{equation}
    Z = \frac{1}{4} n \langle v \rangle,
    \label{eq:flux_isotropic}
\end{equation}
що співпадає з (\ref{eq:flux_formula}), чим підтверджується її загальність.

%% --------------------------------------------------------
\section{133}
%% --------------------------------------------------------

Задача про розподіл молекул газу за швидкістю є чисто класичною в тому сенсі, що швидкості молекул розподілені безперервно і в як завгодно малий інтервал $dv$ «поміщається» досить велика кількість молекул, настільки, що в стаціонарному стані її можна вважати незмінною (для кожного значення $v$).

Тим часом, за певних обставин ці припущення суперечать дослідним даним. Спектр енергії будь-якого фінітного руху, обмеженого малим об’ємом, стає дискретним. Саме ж поняття миттєвої швидкості втрачає сенс.

Принцип невизначеності Гейзенберга:
%
\begin{equation}
    \Delta x \cdot \Delta p \geq \frac{\hbar}{2},
    \label{eq:heisenberg}
\end{equation}
вносить невизначеність в імпульс частинки, а отже і швидкості.

Якщо ці невизначеності порівнянні з тепловими швидкостями руху молекул, то рух буде носити квантовий характер і класичний опис --- неприйнятний. Умову цю можна записати так:
%
\begin{equation}
    \frac{\hbar}{m v} \gtrsim \frac{1}{n^{1/3}}.
    \label{eq:quantum_condition}
\end{equation}

Тут ми вважаємо, що в об’ємі $\Delta V$ перебуває одна частинка, і, отже, $n = 1/V$ є концентрацією молекул газу. Умова класичного руху означає, що ліва частина в (\ref{eq:quantum_condition}) набагато менша правої.

Звідси, виразивши $v$ через температуру, отримуємо, що рух можна вважати класичним (що підкоряється законам Ньютона), і його опис за допомогою максвелловського розподілу прийнятний, якщо:
%
\begin{equation}
    T \gg T_g = \left( \frac{h^2}{2\pi m k} \right)^{3/2} n^{2/3}.
    \label{eq:classical_condition}
\end{equation}

Це і є температурна умова класичного руху.

%% --------------------------------------------------------
\section{134}
%% --------------------------------------------------------

Доречно задатися питанням, наскільки єдиним є розподіл Максвелла. Чи завжди система, маючи в початковий момент часу довільний розподіл, повернеться до максвеллівського? І якщо так, то які властивості руху молекул це забезпечують?

Розподіл Максвелла можна виразити через енергії молекул:
%
\begin{equation}
    f(\vec{p}) \propto e^{-E/kT}.
    \label{eq:maxwell_energy}
\end{equation}

Принцип детальної рівноваги полягає у тому, що в стані рівноваги ймовірність переходу з одного стану в інший дорівнює ймовірності зворотного переходу. Це означає:
%
\begin{equation}
    f(p_1)f(p_2) = f(p_1')f(p_2'),
    \label{eq:detailed_balance}
\end{equation}
де $p_1', p_2'$ --- імпульси після зіткнення.

Із цього принципу можна вивести розподіл Максвелла, як єдиний, що задовольняє умовам детальної рівноваги та збереження енергії та імпульсу при зіткненнях.

Отже, розподіл Максвелла не просто результат статистичного усереднення, а слідство фундаментальних властивостей взаємодії частинок і термодинамічної рівноваги.

%% --------------------------------------------------------
\section{135}
%% --------------------------------------------------------

На перший погляд може здатися, що для того, щоб рух молекул газу був хаотичним, досить, щоб вираз \eqref{eq:detailed_balance} дорівнював нулю лише в середньому. Однак, в такому випадку в русі молекул могла б спостерігатися регулярна компонента, яку можна було б зареєструвати експериментально.

Розглянемо аналогію з футбольними матчами. Якби всі команди були рівносильні, то кожна з них з однаковою ймовірністю могла б програти, виграти або зіграти внічию з будь-якою іншою командою. У цьому випадку принцип детальної рівноваги виконується.

Якщо ж деякі команди систематично виграють, а інші --- програють, то рух в турнірній таблиці стає регулярним, і принцип детальної рівноваги порушується. Аналогічно, якщо в русі молекул є домінуючий напрямок, то розподіл не є Максвелловим.

Це показує, що \emph{хаос} в русі молекул забезпечується не лише \emph{усередненням}, а й \emph{детальною рівновагою} між усіма процесами в системі.




\Closesolutionfile{answer}